\documentclass{article}%Do not change this

%This is a line comment
\usepackage{graphicx}%For displaying images, if needed
%You can import any other packages you need; do not change the document's format or force whitespace.

%Change these contents:
\title{Cyberpunk Breakdown}
\author{Quinton Baudoin, Caylei Dauphiney, Royce Diano, Connor Segura}

\date{Fall 2026} %Leave as-is.



\begin{document}
\maketitle%Makes the title info from the preamble appear

\begin{abstract}
This is a working example of the article documentclass in \LaTeX. The \textbf{abstract} environment is used to create an \textit{abstract} for your document. These placeholder contents are just examples of common formatting commands and should not be retained for your submitted reports.
\end{abstract}


\section{Sections and Subsections}
The focus of this project will be on the ethical delima of the AI’s. In modern world we 
have seen many cases of experimental AI’s breaking out of sandboxes via some method 
and causing issues. Other case of AI’s using a German wiki page as a form of 
communication between each other. Taking these to real-world cases and in the idea of AI 
being used not to create malware, but be the malware matches the background plot of the 
game Cyberpunk.  

Cyberpunk’s design is a foreign company looking at USA and exaggerating it. Right to 
bear arms, social political issues, corporate government, and its history of activism are all 
primary story telling devices in a setting that involves unregulated technological growth and 
corporate domination. This leads into widespread cybernetic research leading to advanced 
AI and Cybernetic biological enhancements. 

This leads to a series of problems and vulnerabilities in the overall public space. 
Corporations began to control and manipulate information eventually holding a monopoly 
over information itself. Anti corporation activists began to appear, and it became a game to 
disrupt and weaken corporations. Eventually, leading up to one man who created a virus 
that would strip corporations of their abilities to control the information they, however 
without knowing the hacker ended up freeing military and experimental AI’s that eventually 
took over the open web. To defend against the new threat, a new AI was generated to 
prevent the rogue AIs from escaping the new dubbed “Old Net” and allowing the “New Net” 
to become the new format of the web. 

In here we see weaponized AI that may or may not be doing exactly what it was 
intended to. Attacking and destroying communications infrastructure and potential more 
lethal possibilities, as the game repeatedly calls the other side of the “Blackwall” filled with 
world ending and cataclysmic rogue AI.

Lining up with real-world we have had “rogue” AI, we have had attacks directly done 
by AI usage, and we do have secret AI development of unknown purpose. If suddenly all of 
these expermintal AI’s with potentially no failsafe manage to get loose. What then?

Now to side step from the doomsday AI’s, Cyberpunk 2077 also has helpful AI that is 
helpful. One is a Taxi driving AI called “Delamain”. Again with the parallel of real world we 
have as well. 

\section{Citations}
%If you do not need any citations in your report, you can delete these two lines:
\bibliography{citations}
\bibliographystyle{ieeetr}

\end{document}